% ==========================================================================
% Capítulo: Motivación y antecedentes
% Estado: se registra en ../STATUS.md, no aquí.
%
% Debe contener (project-context/formato-anteproyecto.md, "Motivación y
% antecedentes"; extensión máxima 3 páginas):
%   - CONTEXTO: escenario (físico, institucional, académico) donde se
%     enmarca el proyecto; involucrados (stakeholders) y sus condiciones;
%     elementos de modo, tiempo y lugar. Toda afirmación sectorial debe ir
%     soportada con cifras/estudios citados (literatura indexada).
%   - ANTECEDENTES DEL PROBLEMA: hechos y circunstancias que motivan o
%     condicionan el proyecto; por qué se presentó el problema (causas y
%     efectos).
%   - JUSTIFICACIÓN: pertinencia y relevancia, en términos de impactos y
%     contribuciones.
% NO debe contener: juicios personales ("creemos", "pensamos que"),
% afirmaciones sin evidencia verificable, ni la formulación del problema
% (esa va en 02-descripcion-problema.tex).
%
% Hechos del proyecto: ../../project-context/documentation.md.
% Mecánica de párrafo: ../../skills/parragraph-structure/, ../../skills/writting-tools/.
% Citas: ../../skills/reference-writting/ y ../../skills/latex/references/citations-and-figures.md.
% ==========================================================================
\chapter{Motivación y antecedentes}\label{cha:motivacion-antecedentes}

\section{Contexto}

Llevar un modelo de aprendizaje automático a producción exige más que entrenarlo. \citet{sculley-hiddentechnicaldebt-2015} señalan que su código es solo una pequeña fracción de un sistema real, rodeado de una infraestructura extensa que incluye el servicio y el monitoreo de los modelos. \citet{kreuzberger-mlopsoverview-2023} definen las operaciones de aprendizaje automático (MLOps) como un paradigma que abarca la conceptualización, implementación, monitoreo, despliegue y escalabilidad de estos productos, e incluyen un componente dedicado a vigilar de forma continua el desempeño en servicio y el estado de la infraestructura. Aun así, \citet{eken-mlopsmultivocalreview-2026} advierten que un cuerpo de conocimiento integrado sobre MLOps sigue sin lograrse, porque su alcance es amplio y abarca desafíos muy diversos de esa transición.

La infraestructura mencionada exige, además, una capacidad de procesamiento creciente. Desde la llegada del aprendizaje profundo a comienzos de la década de 2010, el cómputo necesario para entrenar los sistemas más avanzados se ha duplicado aproximadamente cada seis meses \citep{sevilla-computetrends-2022}. La industria domina una porción cada vez mayor de ese poder de cálculo, mientras la academia cuenta con menos del que necesita para investigar, sobre todo en unidades de procesamiento gráfico (GPU) \citep{ahmed-industryinfluence-2023}. En ese escenario, las GPU que ya posee una universidad adquieren un valor estratégico.

Cuando ese \emph{hardware} se comparte, optimizar su uso se vuelve un reto estructural. Incluso en plataformas empresariales multiusuario de aprendizaje profundo, algunos trabajos presentan una utilización notablemente baja de las GPU asignadas, lo que desperdicia recursos valiosos y reduce la productividad del desarrollo \citep{gao-lowgpuutilization-2024}. El estudio analizó 400 casos reales en un entorno corporativo de gran tamaño, elegidos entre los que aprovechaban en promedio la mitad o menos de su capacidad, e identificó en ellos 706 problemas de este tipo, originados en el código de cada trabajo por cómputo insuficiente en el acelerador o por interrupciones de tareas ajenas a él. Si el problema aparece a esa escala, es razonable que un laboratorio académico también necesite herramientas para detectarlo y controlarlo.

Ese reto se presenta, en una dimensión más modesta, en los equipos de cómputo que el Laboratorio de Inteligencia Artificial y Arquitectura de Software (IAsLab) de la Universidad Icesi destina a la ciencia de datos y la inteligencia artificial. Este proyecto de grado se enmarca en el macroproyecto \emph{Ecosistema computacional de servicios de software para el desarrollo de proyectos del IAsLab para Industria 4.0 y e-Salud en el contexto de la Transformación Digital}, cuyo propósito es reducir la brecha entre la investigación aplicada del IAsLab y su adopción en entornos productivos. Para ello, la iniciativa busca consolidar una base común de componentes y servicios de \emph{software} reutilizables y estandarizados que eleven la madurez tecnológica de sus desarrollos y faciliten su paso de la experimentación académica a soluciones robustas y seguras, acordes con las exigencias de los sectores industrial y clínico.

En ese marco, un proyecto de grado anterior construyó un sistema web que actúa como plano de control (\emph{control plane}) del entrenamiento distribuido. Esa herramienta centraliza la configuración, ejecución y seguimiento de los entrenamientos, aprovisiona los nodos de procesamiento e integra SAAMFI como proveedor institucional de identidad. La presente propuesta aporta al propósito común la etapa siguiente del ciclo, pues construye la plataforma que despliega y sirve los modelos resultantes, gobierna el uso compartido de las GPU y observa el estado del \emph{hardware}, sobre esos mismos nodos.

Alrededor de esa infraestructura intervienen actores con necesidades distintas. Los estudiantes y profesores de la universidad requieren GPU para sus proyectos e investigaciones, y se identifican mediante sus roles institucionales. Los administradores del IAsLab operan los equipos y asumen su soporte. Quienes usan el plano de control anterior comparten esas máquinas, que además deben seguir disponibles para las clases.

\section{Antecedentes del problema}

Antes del proyecto de grado previo, el IAsLab carecía de una plataforma estandarizada para administrar el acceso a su infraestructura de cómputo. La habilitación de usuarios se realizaba mediante credenciales de VPN y permisos otorgados a solicitud individual, lo que permitía el acceso a estudiantes de asignaturas afines como Ingeniería de Software IV, a integrantes del laboratorio y, en general, a docentes o estudiantes que tramitaran su requerimiento. No obstante, dicho trámite dependía de gestiones manuales y en ocasiones conllevaba tiempos de espera prolongados para su aprobación y configuración. La ausencia de un mecanismo automatizado de aprovisionamiento desincentivaba el aprovechamiento continuo de los recursos, por lo cual parte de la comunidad estudiantil recurría a servicios externos como Google Colab o prescindía de cómputo especializado, lo que produjo una subutilización sostenida del \emph{hardware}. A la fecha, el laboratorio emplea sus equipos de forma experimental y todavía no en explotación continua.

Poner un modelo en operación también era un proceso manual. Antes de este proyecto, el laboratorio desplegaba modelos ejecutando a mano \emph{scripts} de Python que instalaban los motores de inferencia y sus dependencias directamente sobre el sistema operativo de las estaciones. En ausencia de una plataforma o interfaz de autoservicio que abstrajera la infraestructura subyacente, este procedimiento no podía gestionarse de forma desacoplada del anfitrión. La puesta en servicio exigía operar directamente sobre los nodos con credenciales de acceso a las máquinas, otorgadas vía VPN y permisos administrativos, condición técnica que supeditaba el despliegue a contar previamente con dicho acceso. El proyecto de grado previo automatizó el aprovisionamiento de nodos para el entrenamiento distribuido, pero no abordó las etapas posteriores al entrenamiento. Hoy ninguna plataforma del laboratorio permite desplegar y servir modelos ya entrenados ni gobernar el uso compartido de las GPU, y los administradores no pueden observar en tiempo real el consumo exacto de la infraestructura.

Las universidades que comparten clústeres de GPU enfrentan problemas semejantes. \citet{xu-sing-2025}, a partir de la operación de un clúster de más de 160 GPU al servicio de más de 480 usuarios de su universidad, señalan que compartir estos recursos es esencial para aprovecharlos y ampliar el acceso, aunque gestionarlos plantea desafíos que van desde la configuración del sistema hasta el reparto equitativo entre usuarios, con personal limitado para operarlos. Según los autores, muchas universidades adoptan una herramienta conocida con configuraciones subóptimas y terminan subutilizando recursos ya limitados. Un esquema simple y común, que entrega a los usuarios acceso directo por consola a las máquinas con GPU, carece de mecanismos de aislamiento y de gestión, obliga a cada usuario a configurar su entorno y a coordinar a mano el uso de las GPU, y dificulta que el reparto sea equitativo.

En una plataforma académica de mayor escala, un mecanismo de gobernanza tuvo efectos medibles. En la National Research Platform, un clúster Kubernetes multiusuario que en 2024 atendió a investigadores de 50 campus, la introducción de un sistema de reservas para las GPU de mayor demanda elevó la utilización promedio de las GPU del clúster del 21.49\,\% al 31.37\,\%, aunque sus operadores reconocen que aún hay margen de mejora \citep{weitzel-nrpreservation-2025}. Estos antecedentes coinciden con la situación del IAsLab, donde el acceso directo y manual a los equipos no ofrece un reparto controlado de las GPU ni visibilidad sobre su uso, condiciones que exponen la infraestructura a un riesgo de monopolización a medida que crece el número de usuarios.

\section{Justificación}

Construir esta plataforma es pertinente porque responde a necesidades concretas del IAsLab. El laboratorio despliega hoy sus modelos de forma manual, no cuenta con reglas que repartan sus GPU entre los usuarios y no dispone de una vista del consumo de su infraestructura, mientras crece la comunidad que podría aprovecharla. Una solución que automatice el despliegue, gobierne el acceso mediante cupos y observe el estado del \emph{hardware} sustituye la entrega manual de credenciales por un servicio al que estudiantes y profesores acceden de forma controlada.

Ese cambio amplía el acceso al cómputo especializado. Al sustituir un esquema de asignación manual y tiempos de espera prolongados por un servicio con reglas de uso explícitas, la capacidad instalada queda al alcance de más estudiantes y profesores, que dejan de depender de suscripciones comerciales que no todos pueden costear, y se fortalecen otros proyectos que requieren esos recursos. La universidad, por su parte, obtiene un mayor retorno de la inversión que ya hizo en equipos, al usarlos en lugar de adquirir créditos en nubes públicas.

Ese mejor aprovechamiento tiene, además, un efecto ambiental. En un clúster universitario compartido, el consumo base de los equipos, que incluye la potencia en reposo de sus componentes, se mantiene relativamente constante sin importar cuántas GPU estén ocupadas, y la GPU representa la mayor parte de la potencia demandada \citep{xu-sing-2025}. Un acelerador asignado que no se usa gasta energía sin producir resultados, de modo que gobernar su reparto y detectar la subutilización favorece un uso más eficiente de la energía del laboratorio.

Por último, el proyecto es relevante porque su escenario reúne características que lo distinguen de los entornos industriales de gran escala. Un laboratorio universitario con pocos nodos de GPU usa las mismas máquinas para las clases y para la investigación, y necesita gobernar el servicio de modelos y observar el estado del \emph{hardware} sobre infraestructura propia. Documentar el diseño, la operación y la evaluación de una plataforma en esas condiciones genera una experiencia transferible a otras instituciones académicas con recursos limitados.
